\section{Local Gromov--Witten evaluation and reconstruction}\label{sec:evaluation}

Let $u$ be the Gromov--Witten variable and $q$ the local curve-degree
variable. We use disconnected maps with no wholly contracted
connected components. Contracted components attached to positive-degree
maps are included. For each partition $\lambda$ let $h(b)$ denote the
hook length of a box $b$ of its Young diagram. The empty partition
contributes $1$ to the following series:
\begin{equation}\label{eq:hook}
 E_g(q,u)=\sum_{\lambda} q^{|\lambda|}
       \prod_{b\in\lambda}\bigl(2\sin(h(b)u/2)\bigr)^{2g-2}.
\end{equation}

\begin{proposition}\label{prop:evaluation}
The disconnected Gromov--Witten evaluation of
$\sum_{m\ge0}a_m(C,L)q^m$ is $E_g(q,u)$.
\end{proposition}
\begin{proof}
Fix $m>0$ and the Euler characteristic of the source.
The stable-map cycle map is proper at fixed Euler characteristic,
and its equivariant virtual pushforward defines the Gromov--Witten
class on $Z_m$ \cite[\S3.5, (3.21)--(3.23)]{Pardon}.
By Lemma~\ref{lem:support} its fixed-locus localization is a class
$v$ at $m[C]$. More precisely, equivariant homology after inverting
$h$ identifies with that of the fixed locus by
\cite[Proposition~4.5]{Pardon}, and $v$ is the inverse image under
fixed-locus pushforward. Indeed, off $m[C]$ the tuple of
Lemma~\ref{lem:support} is a nowhere-zero section. Its Euler
class $(\prod_i r_i)h^n$ vanishes there, so $h^n$ annihilates
the equivariant homology of the complement. The fixed-locus
isomorphism follows from the localization sequence. The class $v$
is a scalar because the virtual dimension is zero.

We check that the compact class \eqref{eq:cut} makes precisely this
scalar an ordinary universal evaluation. Fix a tuple of length $n$
as in Lemma~\ref{lem:support}. At each Borel stage its compact class
restricts on the fixed section to $h^n$. Thus capping with the
equivariant virtual class and pushing forward to the point gives
$h^n v$ after inverting $h$, by the projection formula. The compact
support of the cut-down already defines this pushforward over
$\QQ[h]$, which injects into $\QQ[h,h^{-1}]$. At stage $N=n$ the
ordinary integral is $v$ since $\int_{B_n}h^n=1$. This is the
evaluation of $a_m(C,L)$, by the finite Borel description in
\cite[Proposition~4.8]{Pardon}.

It remains to identify $v$. A fixed map has fixed image cycle, so
all its positive-degree components map into $C$. Attached contracted
components map into $C$ as well. Conversely every stable map into
$C$ is fixed. Write $\mathcal F$ for this proper fixed moduli stack
and $p$ for the cycle map. Along $C$ the tangent bundle splits as
$T_C\oplus L_+\oplus L_-$. The fixed part of the obstruction theory
is the stable-map theory of $C$; the moving virtual normal complex
is $R\pi_*f^*L_+\oplus R\pi_*f^*L_-$ with the respective
characters $+1,-1$, where $(\pi,f)$ is the universal curve and map.
Both characters stay nonzero on the antidiagonal. Equivariant
virtual localization and proper pushforward give
\[
 v=(p|_{\mathcal F})_*
    \left(\frac{[\mathcal F]^{\mathrm{vir}}}
    {e_{S^1}(N^{\mathrm{vir}})}\right).
\]
This is the residue integral of \cite[\S2.2]{BP}. Its degree-$m$
normalization shifts the power of $u$ by
$m(2-2g+\deg L_++\deg L_-)=0$. Formula~\eqref{eq:hook} now
follows from \cite[Corollary~7.3]{BP}.
\end{proof}

The degree-two term distinguishes the balanced degrees among all
splittings with determinant $K_C$. Indeed, take line degrees
$g-1+\ell,g-1-\ell$ and retain the antidiagonal action. Write
$E_{g,\ell}$ for its disconnected residue series, so $E_{g,0}=E_g$.
In the formula of \cite[Corollary~7.2]{BP} the change of degrees in
degree two multiplies the two partition contributions by
$e^{i\ell u}$ and $e^{-i\ell u}$. Indeed, the partitions $(2)$ and
$(1,1)$ have opposite total contents $\pm1$ and the same hook
lengths, and the sign in degree two is $(-1)^{2\ell}=1$. Thus
\begin{equation}\label{eq:balance}
 [q^2]E_{g,\ell}=\cos(\ell u)\,[q^2]E_g .
\end{equation}
The balanced coefficient is a nonzero Laurent series. Equality
with the elementary-cluster coefficient therefore forces $\ell=0$.

To pass from this evaluation to a class in $H$, we extend the
reconstruction argument of \cite[\S2]{Transition}. Write
$x_{h,j}=x_{h,j,0}$ for Pardon's scalar generators,
$x_{h,0}=1$, and set
\[
 f_h(q)=\sum_{j\geq0}x_{h,j}q^j,\qquad
 p_{h,j}=[q^j]\log f_h(q)\quad(j>0).
\]
Over $\QQ$ the $p_{h,j}$ are polynomial coordinates on $H$.
Pardon's scalar formulas \cite[Lemmas~4.12--4.13 and
Theorem~7.6]{Pardon} give \eqref{eq:operations} for $x_{h,m}$.
Thus $p_{h,j}$ is primitive for $\Delta$, and $\rho_d p_{h,j}$
is $p_{h,j/d}$ if $d\mid j$ and zero otherwise.

Let $\Gamma$ be a submonoid of a finitely generated abelian group with
an additive degree on $\Gamma$, positive and integral off zero, and
with finite sublevel sets. Complete the monoid ring in this degree.
Call a series
$U=\sum_\beta U_\beta Q^\beta$ \emph{compatible} if $U_0=1$, its
other coefficients lie in the augmentation ideal of $H$, and
\begin{equation}\label{eq:compat}
 \Delta U_\beta=\sum_{\gamma+\eta=\beta}U_\gamma\otimes U_\eta,
 \qquad \rho_d U_\beta=\sum_{\substack{\gamma\in\Gamma\\ d\gamma=\beta}}U_\gamma .
\end{equation}
The second identity says that $\rho_d U=\psi_d U$, where
$\rho_d$ acts on coefficients and $\psi_d$ substitutes
$Q^\gamma\mapsto Q^{d\gamma}$. The latter is a continuous ring
endomorphism even when multiplication by $d$ identifies distinct
classes. The coefficient at $\beta$ sums all their contributions.

\begin{lemma}\label{lem:reconstruction}
A compatible series is determined by its disconnected
Gromov--Witten evaluation. It admits a unique product expansion
in the $f_h(Q^\beta)$ with rational exponents and finite genus
support in each class.
\end{lemma}
\begin{proof}
The primitive polynomials in the $p_{h,j}$ are precisely their
finite linear combinations, since an additive polynomial over
$\QQ$ is linear. If a primitive $R$ is killed by every $\rho_d$
for $d>1$, it is a linear combination of $p_{h,1}=x_{h,1}$.
Indeed applying $\rho_j$ for the largest multiplicity $j>1$ in
such a polynomial detects all its coefficients at $j$. Set
$w=(2\sin(u/2))^2$. The scalar degree-one evaluations are
$\GW(x_{h,1})=w^{h-1}$ by the degree-one local formula
\cite[\S8]{BP}. These functions are linearly independent by their
distinct lowest powers of $u$. If two compatible series differ,
choose a class $\beta$ of least degree at which they differ.
Every decomposition of $\beta$ into two nonzero summands uses
lower degrees. The first identity in \eqref{eq:compat} therefore
makes the difference at $\beta$ primitive. For $d>1$ each root
$d\gamma=\beta$ has degree $\deg(\beta)/d$, so the second identity
kills the difference under $\rho_d$. Equal evaluations force it
to vanish.

For the expansion, every coefficient of $\log U$ is primitive.
Indeed the coproduct identity gives $\Delta U=(U\otimes1)(1\otimes U)$
with one common monoid variable, so taking the logarithm gives the
primitive identity coefficientwise. Write
$[Q^\beta]\log U=\sum_{h,j}b_{\beta;h,j}p_{h,j}$. Since $\rho_j$
and $\psi_j$ commute with logarithms, comparison of the $p_{h,1}$
coefficients gives
\[
 b_{\beta;h,j}=\sum_{j\gamma=\beta} b_{\gamma;h,1}.
\]
This is exactly the $p_{h,j}Q^\beta$ coefficient of
$\sum_{\gamma\ne0,h}b_{\gamma;h,1}\log f_h(Q^\gamma)$.
The sums are finite in each degree, because positive-class degrees
are at least one and each coefficient in $H$ is a polynomial.
Exponentiating gives the expansion. Its exponents are unique since
$p_{h,1}Q^\gamma$ in the logarithm recovers $b_{\gamma;h,1}$.
\end{proof}

\begin{corollary}\label{cor:characterization}
Assume the universal Gromov--Witten evaluation for analytic families.
For each $g$ there is exactly one compatible series in $H\llbracket q\rrbracket$
whose disconnected evaluation is the contribution of a genus-$g$
elementary cluster. Every balanced antidiagonal model of genus $g$
represents this series.
\end{corollary}
\begin{proof}
Propositions~\ref{prop:operations} and~\ref{prop:evaluation} give
existence, since the elementary-cluster contribution is $E_g$
\cite[Proposition~3.3]{IP}. Lemma~\ref{lem:reconstruction} with
$\Gamma=\mathbb N$ gives uniqueness. In particular the coefficients
are independent of the curve, the line bundle, and the order of
the two summands. We denote them by $a_{g,m}$ and the series by $A_g$.
\end{proof}
